\heading{实验物理中的统计方法-模拟题1-期末}
\noteinfo{题目和答案由AI生成。知识点范围按现有往年题整理，叙述风格尽量向往年题靠拢；本套难度较往年期末略高。}

\begin{question}
设物理量
\[
Z=\frac{X}{Y},
\]
其中 $X,Y$ 的标准差分别为 $\sigma_X,\sigma_Y$，协方差为 $c=\cov(X,Y)$。写出小误差近似下 $Z$ 的不确定度公式。
\begin{solution}
由误差传播公式，
\[
\sigma_Z^2\approx
\left(\frac{\partial Z}{\partial X}\right)^2\sigma_X^2+
\left(\frac{\partial Z}{\partial Y}\right)^2\sigma_Y^2+
2\frac{\partial Z}{\partial X}\frac{\partial Z}{\partial Y}c.
\]
而
\[
\frac{\partial Z}{\partial X}=\frac1Y,\qquad
\frac{\partial Z}{\partial Y}=-\frac{X}{Y^2}.
\]
故
\end{solution}
\end{question}

\begin{question}
设 $X$ 服从自由度为 $2n$ 的卡方分布。$n$ 很大时，$\sqrt{X}$ 近似服从什么分布？
\begin{solution}
对卡方分布，
\[
E[X]=2n,\qquad \V(X)=4n.
\]
令
\[
g(x)=\sqrt{x},
\qquad g'(x)=\frac{1}{2\sqrt{x}}.
\]
由 delta 方法，
\[
\sqrt{X}\approx N\!\left(g(2n),\ [g'(2n)]^2\V(X)\right).
\]
因此
\[
[g'(2n)]^2\V(X)
=\frac{1}{4\cdot 2n}\cdot 4n=\frac12.
\]
故
\end{solution}
\end{question}

\begin{question}
设总体的概率密度为
\[
f(x;\theta)=\begin{cases}
(\theta+1)x^\theta,&0<x<1,\\
0,&\text{其他},
\end{cases}
\qquad \theta>-1.
\]
现有样本 $X_1,\cdots,X_n$。
\begin{enumerate}
\item 求 $\theta$ 的矩估计量；
\item 求 $\theta$ 的极大似然估计量；
\item 求极大似然估计量的渐近方差。
\end{enumerate}
\begin{solution}
先算总体均值：
\[
E[X]=\int_0^1 x(\theta+1)x^\theta\,dx
=\frac{\theta+1}{\theta+2}.
\]
令 $\bar X=E[X]$，解得矩估计
\[
\hat\theta_{\rm MOM}=\frac{2\bar X-1}{1-\bar X}.
\]
似然函数为
\[
L(\theta)=\prod_{i=1}^n (\theta+1)x_i^\theta,
\qquad
\ell(\theta)=n\ln(\theta+1)+\theta\sum_{i=1}^n\ln x_i.
\]
求导并令零，
\[
\frac{n}{\theta+1}+\sum_{i=1}^n \ln x_i=0,
\]
故
\ansmath{\hat\theta_{\rm MLE}=-\frac{n}{\sum_{i=1}^n\ln x_i}-1}
Fisher 信息为
\[
I_1(\theta)=-E\!\left[\frac{\partial^2 \ell_1}{\partial \theta^2}\right]
=\frac{1}{(\theta+1)^2}.
\]
所以
\[
I_n(\theta)=\frac{n}{(\theta+1)^2}.
\]
于是 MLE 的渐近方差为
\ansmath{\V(\hat\theta_{\rm MLE})\approx \frac{(\theta+1)^2}{n}}
\end{solution}
\end{question}

\begin{question}
设 $X_1,\cdots,X_n$ 独立同分布，且都服从参数为 $\lambda$ 的泊松分布。求 $\lambda$ 的极大似然估计，并说明它是有效估计量。
\begin{solution}
似然函数为
\[
L(\lambda)=\prod_{i=1}^n e^{-\lambda}\frac{\lambda^{x_i}}{x_i!},
\qquad
\ell(\lambda)=-n\lambda+\Bigl(\sum_{i=1}^n x_i\Bigr)\ln\lambda+\text{const}.
\]
令导数为零，
\[
-n+\frac{\sum x_i}{\lambda}=0,
\]
得
\[
\hat\lambda=\bar X.
\]
又有
\[
E[\bar X]=\lambda,\qquad \V(\bar X)=\frac{\lambda}{n}.
\]
单个样本的 Fisher 信息为
\[
I_1(\lambda)=\frac{1}{\lambda},
\]
因此
\[
I_n(\lambda)=\frac{n}{\lambda},
\qquad
\frac{1}{I_n(\lambda)}=\frac{\lambda}{n}.
\]
由于
\[
\V(\hat\lambda)=\frac{\lambda}{n}
\]
恰好达到 Cram\'er--Rao 下界，因此
\anstext{$\hat\lambda=\bar X$ 是有效估计量。}
\end{solution}
\end{question}

\begin{question}
从某总体中抽取 $16$ 个样本，得到
\[
\bar X=10.4,\qquad S=1.2.
\]
假设总体服从正态分布，检验
\[
H_0:\mu=10
\]
是否能够在 $95\%$ 置信水平下被接受，并给出 $\mu$ 的 $95\%$ 置信区间。已知
\[
t_{0.975}(15)=2.131.
\]
\begin{solution}
作双侧 $t$ 检验。统计量为
\[
T=\frac{\bar X-\mu_0}{S/\sqrt{n}}
=\frac{10.4-10}{1.2/4}
=\frac{0.4}{0.3}
\approx 1.333.
\]
由于
\[
|T|=1.333<2.131,
\]
故不能拒绝 $H_0$，也就是说在 $95\%$ 置信水平下，可以认为数据与 $\mu=10$ 相容。
再求置信区间：
\[
\bar X\pm t_{0.975}(15)\frac{S}{\sqrt{n}}
=10.4\pm 2.131\times 0.3
=10.4\pm 0.6393.
\]
因此
\ansmath{\mu\in (9.761,\ 11.039)\qquad (95\%\ \text{CI})}
\end{solution}
\end{question}

\begin{question}
设 $X_1,\cdots,X_n$ 独立同分布于 $U(0,\theta)$，记
\[
Y=\max(X_1,\cdots,X_n).
\]
\begin{enumerate}
\item 求 $\theta$ 的极大似然估计量；
\item 求一个由 $Y$ 构造的无偏估计量；
\item 比较二者的均方误差。
\end{enumerate}
\begin{solution}
由似然函数
\[
L(\theta)=\theta^{-n}\mathbf 1(\theta\ge Y),
\]
可知
\[
\hat\theta_{\rm MLE}=Y.
\]
又因为
\[
E[Y]=\frac{n}{n+1}\theta,
\]
所以无偏估计量可取
\[
\tilde\theta=\frac{n+1}{n}Y.
\]
再算
\[
\V(Y)=\frac{n\theta^2}{(n+1)^2(n+2)},
\qquad
\Bias(Y)=-\frac{\theta}{n+1}.
\]
因此
\[
\MSE(Y)=\V(Y)+\Bias(Y)^2
=\frac{2\theta^2}{(n+1)(n+2)}.
\]
而
\[
\MSE(\tilde\theta)=\V(\tilde\theta)
=\left(\frac{n+1}{n}\right)^2\V(Y)
=\frac{\theta^2}{n(n+2)}.
\]
比较得
\[
\frac{1}{n(n+2)}<\frac{2}{(n+1)(n+2)}\qquad (n>1).
\]
故
\anstext{当 $n>1$ 时，无偏估计量 $\tilde\theta$ 的均方误差更小。}
\end{solution}
\end{question}

\begin{question}
测量模型为
\[
y_i=ax_i+\varepsilon_i,\qquad i=1,\cdots,n,
\]
其中 $x_i$ 已知，$\varepsilon_i$ 相互独立且
\[
\varepsilon_i\sim N(0,\sigma_i^2).
\]
求参数 $a$ 的最小二乘估计、该估计量的方差，以及 $\chi^2_{\min}$。
\begin{solution}
加权最小二乘的目标函数为
\[
\chi^2(a)=\sum_{i=1}^n \frac{(y_i-ax_i)^2}{\sigma_i^2}.
\]
对 $a$ 求导并令零：
\[
\frac{d\chi^2}{da}=-2\sum_{i=1}^n \frac{x_i(y_i-ax_i)}{\sigma_i^2}=0.
\]
因此
\[
\hat a=\frac{\sum_{i=1}^n \dfrac{x_i y_i}{\sigma_i^2}}
{\sum_{i=1}^n \dfrac{x_i^2}{\sigma_i^2}}.
\]
又由于 $\hat a$ 是 $y_i$ 的线性组合，故
\[
\V(\hat a)=\frac{1}{\sum_{i=1}^n \dfrac{x_i^2}{\sigma_i^2}}.
\]
代回可得
\[
\chi^2_{\min}
=\sum_{i=1}^n \frac{y_i^2}{\sigma_i^2}
-\frac{\left(\sum_{i=1}^n \dfrac{x_i y_i}{\sigma_i^2}\right)^2}
{\sum_{i=1}^n \dfrac{x_i^2}{\sigma_i^2}}.
\]
\end{solution}
\end{question}

\begin{question}
一次稀有事例实验中，信号区计数
\[
N_s\sim \mathrm{Poisson}(s+b),
\]
控制区计数
\[
N_c\sim \mathrm{Poisson}(3b),
\]
且二者独立。若观测到
\[
N_s=6,\qquad N_c=9,
\]
求 $s,b$ 的极大似然估计。
\begin{solution}
联合对数似然为
\[
\ell(s,b)=6\ln(s+b)-(s+b)+9\ln(3b)-3b+\text{const}.
\]
对 $s$ 求偏导：
\[
\frac{\partial \ell}{\partial s}=\frac{6}{s+b}-1=0
\quad\Longrightarrow\quad
s+b=6.
\]
对 $b$ 求偏导：
\[
\frac{\partial \ell}{\partial b}=\frac{6}{s+b}-1+\frac{9}{b}-3=0.
\]
由 $s+b=6$ 可知前两项抵消，于是
\[
\frac{9}{b}-3=0\quad\Longrightarrow\quad b=3.
\]
从而
\[
s=6-b=3.
\]
故
\ansmath{\hat b=3,\qquad \hat s=3}
\end{solution}
\end{question}

